% Section 8 — Discussion and practitioner decision table (compressed)
\section{Discussion and Decision Table}
\label{sec:decision}

\subsection{What the boundary analysis implies for practitioners}
\label{sec:dec-implications}

The boundary analysis answers the three questions as follows. For RQ1, Kafka EOS fails to
preserve ledger invariants exactly when the external write and the progress can be committed in
different orders across a crash or unknown-outcome timeout: the duplicate window (effect durable,
progress not) and the loss window (progress durable, effect not). Broker failure, consumer death,
rebalance and partition are not failures \emph{of Kafka} here; they expose the two windows. For
RQ2, mechanisms differ in \emph{which} window they close and at what cost, not in a single score.
For RQ3, the minimal sufficient set depends on the exposure the operator accepts and on whether
the external effect must be independently recoverable.

The practical consequence is that a deployment claim should name the mechanism, not the label:
``we use Kafka exactly-once'' says nothing about the external store, whereas ``we commit the
ledger and the progress in one database transaction'' (M4) says everything about both windows,
with a residual zombie-writer exposure that an idempotency index closes (M4+M3). Operators who
cannot move progress into the database can use an atomic ledger transaction plus a deduplication
claim (M2), at the cost of an inbox/outbox write and a relay; operators who must keep the Kafka
transaction for the publication stream can add an index to the external write (M1x+M3), which
closes duplicates but not losses.

\subsection{Practitioner decision table}
\label{sec:dec-table}

The decision summary below maps a failure exposure to the minimal mechanism set that closes it, the
model's argument for that coverage, and what the executed reduced programme observed. It is
derived from the formal results of Section~\ref{sec:formal} and the descriptive evidence of
Section~\ref{sec:results}; because the reduced programme has $N=3$ per cell and no measured
fixed-rate/saturation performance block, the cost is stated qualitatively rather than as a
measured contrast. Every entry is either a model-derived claim or an observation traceable to
\texttt{results/analysis-reduced/}; none is set by hand.

\paragraph{Decision summary.} (i) $W_1$ duplicate exposure: the minimal sufficient mechanism is
an idempotency key (M3) or a transactional outbox (M2); M4 includes it atomically. Reduced
evidence: K3 duplicated only in the unguarded M0/M1x, while M1x+M3, M2, M2+M3 and M3 were clean.
(ii) $W_2$ lost-effect exposure: no mechanism in the suite closes the early-ordering probe;
prevention requires effect-before-progress durability plus reconciliation. Reduced evidence: the
K4 probe lost postings in all 24 of 24 runs (272--494 per cell), across every exposed mechanism.
(iii) $K_1$/$K_2$ (nothing durable, or unknown outcome): plain at-least-once suffices when the
effect is retried and either idempotent or absent; both classes were clean in 44 runs, including
M0. (iv) Scale (100k+, saturation): depends on load and is not measured here; the reduced
programme reports committed-event throughput only, with no ranking claim, and the pre-registered
Stage~B/C remain future work on the intended host.

A caution accompanies the table: zero-observed cells carry a rule-of-three upper bound, evidence
of practical adequacy at the tested sample size, not a proof of impossibility. The word
``certified'' is reserved, per the pre-registered language rule, for cells where the formal model
proves the invariant under the mechanism's assumptions \emph{and} the deterministic adversarial
runs are clean; a purely empirical zero is reported as ``no violation observed in $N$ runs''.

\subsection{When the extra machinery is not worth it}
\label{sec:dec-negative}

Not every pipeline needs atomic progress. If duplicates can be corrected upstream before they
affect balances, or the sink is idempotent by construction, an at-least-once pipeline with an
idempotency index may suffice and be cheaper. The decision table therefore begins with the
exposure question --- which window is reachable, what a duplicate or lost posting costs, and
whether the external effect is independently recoverable --- before the measured cost column
matters.

\subsection{Relation to prior work}
\label{sec:dec-prior}

The model is consistent with the vendor's statement of the boundary~\cite{kafka-design} and
KIP-98's producer-scoped guarantee~\cite{kip98}, but makes structural what those documents state
qualitatively: both windows are reachable under ordinary fault schedules. It extends
stream-processing fault-tolerance results~\cite{flink2017,millwheel2013} by showing that
snapshot-style reasoning protects internal state but not an external database, gives the
outbox/idempotency patterns~\cite{debezium-outbox,stripe-idempotency} a correctness--cost profile
in the ledger setting, and supplies the consistency contract missing from microservice
surveys~\cite{laigner2021,micro2026}.

\subsection{Reusability of the method}
\label{sec:dec-reuse}

The method is mechanism-agnostic: a ground-truth generator that writes the expected effect
before producing, an independent oracle that recomputes it, scripted faults with fire receipts,
and exact statistics for mostly-zero outcomes. Any pipeline writing to an external store can be
placed in the same harness by instantiating the poll/commit skeleton and declaring which durable
transitions share a site, and the bounded TLA+ model can be extended with the new mechanism's
transition relation and checked at the same bounds before an expensive campaign.
